
\section{Exact theory of bremsstrahlung in the ultrarelativistic case}

\SourcePage{465}{470}

The matrix element for bremsstrahlung
\begin{equation}
M_{fi} = \int\psi_{\varepsilon'\mathbf{p}'}^{(-)*}(\boldsymbol{\alpha}\mathbf{e}^*)\,e^{-i\mathbf{k}\mathbf{r}}\,\psi_{\varepsilon\mathbf{p}}^{(+)}\,d^3x;
\label{eq:96.1}
\end{equation}
the wave functions of the initial ($\varepsilon$, $\mathbf{p}$) and final ($\varepsilon'$, $\mathbf{p}'$) electrons contain in their asymptotics respectively an incoming and outgoing spherical wave. The computation of this integral is analogous to the computation of the matrix element (95.2). We shall, however, present here a different method of computing the bremsstrahlung cross section, based on the quasi-classicality of the process and not utilising the explicit form of the wave functions of the electron in the field of the nucleus; in this sense the method is not connected with the specific form of the potential of the field (\emph{V.~N.~Baier, V.~M.~Katkov}, 1968).

In the process of bremsstrahlung the nucleus transfers to the electron and photon a momentum $\mathbf{q} = \mathbf{p}' + \mathbf{k} - \mathbf{p}$. As in the problem of pair production, one must distinguish two regions of values of the momentum transfer $q_\perp$, transverse with respect to $\mathbf{p}$:
\begin{equation}
\text{I)}\; m \gtrsim q_\perp \gg \omega m^2/\varepsilon^2, \quad \text{II)}\; q_\perp \sim \omega m^2/\varepsilon^2 \ll m.
\label{eq:96.2}
\end{equation}

It is evident that in region I the cross section of photon emission is given by its Born value: for such $q_\perp$ the change in the recoil momentum of the nucleus upon radiation is inessential, as will be shown in \S\,98 (see the derivation of condition (98.10)). Therefore in region I the cross section of the process equals the product of the exact cross section of scattering of the electron in the field of a fixed nucleus and the probability of photon emission, which is independent of the form of the field. But according to (80.10) the scattering cross section in a Coulomb field for small angles coincides with its Born value. The same therefore applies to the cross section of the whole process in region I.

Thus, only region II requires special consideration. To small momentum transfers there corresponds the passage of the electron past the nucleus at large impact parameters: $\rho \sim 1/q_\perp \gtrsim \varepsilon/m^2$

\SourcePage{466}{471}

$\gtrsim \varepsilon/m^2$. But at such distances the motion of the electron is certainly quasi-classical, as one can easily verify by a simple application of the usual quasi-classicality condition (46.7) (see III) to the ultrarelativistic equation (39.5).

The quasi-classicality of the motion permits applying the method already used in \S\,90 for magnetobremsstrahlung radiation. In this case the expression (90.7) represents the probability of emission during a single passage of the electron past the nucleus.

For the function $L$ that appeared in \S\,90 the formula (90.18) remains in force; the sole distinction consists in the form of the quasi-classical trajectory of the electron $\mathbf{r} = \mathbf{r}(t)$, along which the difference $\mathbf{r}_2 - \mathbf{r}_1$ is computed.

At large impact parameters the field of the nucleus can be considered weak. In the zeroth approximation the trajectory is a straight line passing at a distance $\rho$ from the centre. In the next approximation we have the equation of motion (cf.\ I, \S\,20)
\[
\frac{d\mathbf{p}}{dt} = -\frac{\boldsymbol{\rho}}{r}\,\frac{dU}{dr},
\]
where $\boldsymbol{\rho}$ is a vector in the $xy$ plane, perpendicular to the initial momentum of the electron, and for $r$ on the right-hand side of the equation one should take the zeroth-approximation function:
\[
r \approx \sqrt{\rho^2 + v^2t^2} \approx \sqrt{\rho^2 + t^2}.
\]
Consequently,
\begin{equation}
\mathbf{p}(t) - \mathbf{p}_1 = -\boldsymbol{\rho}\int_{t_1}^{t}\frac{dU}{dr}\,\frac{dt}{r}.
\label{eq:96.3}
\end{equation}
With sufficient accuracy the velocity $\mathbf{v}(t) = \mathbf{p}(t)/\varepsilon$ (where the energy $\varepsilon$ depends only on the magnitude, but not on the direction of $\mathbf{p}$) can be considered constant. One more integration then gives
\begin{equation}
\mathbf{r}(t) - \mathbf{r}_1 = \mathbf{v}_1(t - t_1) - \frac{1}{\varepsilon}\int_{t_1}^{t}[\mathbf{p}(t') - \mathbf{p}_1]\,dt'.
\label{eq:96.4}
\end{equation}

Setting $t_1 = -\infty$, the quantities $\mathbf{p}_1 = \mathbf{p}(-\infty) \equiv \mathbf{p}$ and $\mathbf{v} = \mathbf{p}/\varepsilon$ will be the initial momentum and velocity of the electron.

Let us represent the probability (90.7) in the form
\begin{equation}
dw = |a(\boldsymbol{\rho})|^2\,\frac{d^3k}{(2\pi)^3},
\label{eq:96.5}
\end{equation}
where
\begin{equation}
a(\boldsymbol{\rho}) = e\sqrt{\frac{2\pi}{\omega}}\int_{-\infty}^{\infty}R(t)\exp\!\left\{i\frac{\varepsilon}{\varepsilon'}[\omega t - \mathbf{k}\mathbf{r}(t)]\right\}dt,
\label{eq:96.6}
\end{equation}
\[
R(t) = \frac{u_{\varepsilon'\mathbf{p}'}^*}{\sqrt{2\varepsilon'}}\,(\boldsymbol{\alpha}\mathbf{e}^*)\,\frac{u_{\varepsilon\mathbf{p}}}{\sqrt{2\varepsilon}},
\]

\SourcePage{467}{472}

where $\varepsilon' = \varepsilon - \omega$, $\mathbf{p}'(t) = \mathbf{p}(t) - \mathbf{k}$. The classical function $\mathbf{p}(t)$ is given by formula (96.3). If $\mathbf{p}$ is the initial momentum of the electron, then for a Coulomb field ($U = -\nu/r$, $\nu = Z\alpha$) we have
\[
\mathbf{p}(t) = \mathbf{p} - \frac{\nu\boldsymbol{\rho}}{\rho^2}\!\left[\frac{t}{\sqrt{\rho^2 + t^2}} + 1\right],
\]
\[
\mathbf{r}(t) = \frac{\mathbf{p}}{\varepsilon}\,t - \frac{\boldsymbol{\rho}}{\rho^2}\,\frac{\nu}{\varepsilon}\!\left[\sqrt{\rho^2 + t^2} + t\right].
\]

Introducing the momentum transfer in classical scattering
\begin{equation}
\boldsymbol{\Delta} = \mathbf{p}(\infty) - \mathbf{p}(-\infty) = -\frac{2\boldsymbol{\rho}\nu}{\rho^2},
\label{eq:96.7}
\end{equation}
we can rewrite these formulas as
\begin{equation}
\begin{split}
\mathbf{p}(t) &= \mathbf{p} + \frac{1}{2}\boldsymbol{\Delta}\!\left[\frac{t}{\sqrt{\rho^2 + t^2}} + 1\right], \\
\mathbf{r}(t) &= \left(\mathbf{p} + \frac{1}{2}\boldsymbol{\Delta}\right)\frac{t}{\varepsilon} + \frac{\boldsymbol{\Delta}}{2\varepsilon}\sqrt{t^2 + \rho^2}.
\end{split}
\label{eq:96.8}
\end{equation}

Using now formula (90.20) for $R(t)$ and the expressions (96.8) for $\mathbf{p}(t)$ and $\mathbf{r}(t)$, one can carry out the integration over time in (96.6). It is accomplished by the introduction of the variable
\[
\xi = -\frac{\varepsilon}{\varepsilon'}[\omega t - \mathbf{k}\mathbf{r}(t)]
\]
instead of $t$ and the use of the formula
\[
\int_{-\infty}^{\infty}\frac{\xi\,e^{-i\xi}\,d\xi}{\sqrt{\chi^2 + \xi^2}} = 2i\chi K_1(\chi),
\]
where $K_1$ is the MacDonald function. In the full execution of this computation, however, there is no necessity, since we require the expression $a(\boldsymbol{\rho})$ only for small values of the independent parameter $\Delta$ ($\Delta \ll m$). In this case we find
\begin{equation}
a(\boldsymbol{\rho}) = w_f^*\mathbf{D}w_i\cdot\Delta\chi K_1(\chi),
\label{eq:96.9}
\end{equation}
where
\[
\chi = \rho\,\frac{\omega\varepsilon}{\varepsilon'}(1 - \mathbf{n}\mathbf{v}),
\]
$\mathbf{n} = \mathbf{k}/\omega$, and $\mathbf{D}$ is some function of $\mathbf{p}$, $\varepsilon$ and $\mathbf{k}$ (but not $\boldsymbol{\rho}$), the precise form of which is inessential\footnote{The spinors $w_i$ and $w_f$ can be considered constant during the integration, i.e.\ one can neglect the change of polarisation of the electron during its classical ultrarelativistic motion. This follows from the equations obtained in \S\,41.}.

\SourcePage{468}{473}

Since in the ultrarelativistic case the photon is emitted at a small angle $\theta$ to the direction of velocity of the electron, we have
\[
\chi \approx \rho\,\frac{\varepsilon}{\varepsilon'}\,\omega\!\left(1 - v + \frac{\theta^2}{2}\right),
\]
or
\begin{equation}
\chi = \rho\,\frac{\omega m^2}{2\varepsilon\varepsilon'}(1 + \delta^2), \quad \delta = \frac{\theta\varepsilon}{m}.
\label{eq:96.10}
\end{equation}

It was already mentioned that (96.5) is the probability of photon emission during a single passage of the electron past the nucleus at impact parameter $\rho$. The cross section of emission of a photon with given frequency and direction is obtained by multiplying this probability by $dp_x\,dp_y/v \approx dp_x\,dp_y \equiv d^2\rho$ and integrating over the impact parameters:
\begin{equation}
d\sigma = \frac{d^3k}{(2\pi)^3}\int|a(\boldsymbol{\rho})|^2\,d^2\rho.
\label{eq:96.11}
\end{equation}

One should not, however, think that this formula without the integration over $d^2\rho$ would also give the distribution of the final electrons in direction. The deflection of the electron during its motion along the classical orbit is determined uniquely by the external field and certainly does not coincide with the uncertain quantum-mechanical deflection (and the limiting value $\mathbf{p}'(\infty)$ of the classical function $\mathbf{p}'(t)$ does not therefore coincide with the real values of the momentum of the electron). To find this distribution it is necessary, therefore, to expand the wave function of the electrons in plane waves.

As is evident from (96.11), $a(\boldsymbol{\rho})$ is the amplitude of photon emission in a collision at impact parameter $\boldsymbol{\rho}$. From the expressions (96.5), (96.6) this amplitude is determined only up to a phase factor. The latter is, obviously, $e^{-i\mathbf{k}\boldsymbol{\rho}}$---by virtue of the presence of the time-independent term $r_\perp(\infty) = \boldsymbol{\rho}$ in $\mathbf{r}(t)$, this constant factor must appear in $V_{fi}(t)$ and can be taken outside the integral sign. Since it is not an operator, it is not affected by commutation operations and, thus, the amplitude of the emission process is
\begin{equation}
e^{-i\mathbf{k}\boldsymbol{\rho}}\,a(\boldsymbol{\rho}),
\label{eq:96.12}
\end{equation}
where $a(\boldsymbol{\rho})$ is given by expression (96.9).

Let the electron be described as $z \to -\infty$ by a plane wave with momentum $\mathbf{p}$, directed along the $z$ axis. This means that the wave function of the electron as $z \to -\infty$ does not depend on $x$ and $y$ and reduces to a constant, which can be set equal to 1.

\SourcePage{469}{474}

Then the wave function of the electron that has passed through the field, as $z \to \infty$, equals\footnote{Cf.\ III, (131.4). We have in mind here the analogy between equation (39.5) (in which we set $p^2 \approx \varepsilon^2$) and the non-relativistic Schr\"{o}dinger equation. Taking into account the difference of coefficients in these equations, it is easy to see that in our case the condition (131.1) (see III) of applicability of equation (131.4) (see III) is indeed satisfied. The fact that this formula does not apply to the region of arbitrarily large $z$ is inessential for the same reasons as in III, \S\,131.}
\begin{equation}
\psi(\infty) = S(\boldsymbol{\rho}) = \exp\!\left\{-i\int_{-\infty}^{\infty}U(x,y,z)\,dz\right\}.
\label{eq:96.13}
\end{equation}

On the other hand, by the meaning of the transition amplitude (96.12) the wave function of the electron that has passed through the field and emitted a photon is
\begin{equation}
e^{-i\mathbf{k}\boldsymbol{\rho}}\,a(\boldsymbol{\rho})\,S(\boldsymbol{\rho}).
\label{eq:96.14}
\end{equation}

The amplitude of the process of photon emission, in which the electron remains in a state with definite momentum $\mathbf{p}'$, is given by the corresponding Fourier component of the function (96.14), i.e.
\begin{equation}
a(\mathbf{q}_\perp) = \int e^{-i\mathbf{p}'\boldsymbol{\rho}}\,e^{-i\mathbf{k}\boldsymbol{\rho}}\,a(\boldsymbol{\rho})\,S(\boldsymbol{\rho})\,d^2\rho = \int e^{-i\mathbf{q}_\perp\boldsymbol{\rho}}\,a(\boldsymbol{\rho})\,S(\boldsymbol{\rho})\,d^2\rho,
\label{eq:96.15}
\end{equation}
where $\mathbf{q}_\perp$ is the transverse component of the momentum transfer vector to the nucleus (cf.\ III, (131.7)). The cross section of scattering with given value of $\mathbf{q}_\perp$ is then
\begin{equation}
d\sigma = |a(\mathbf{q}_\perp)|^2\,\frac{d^3k}{(2\pi)^3}\,\frac{d^2q_\perp}{(2\pi)^2}.
\label{eq:96.16}
\end{equation}

Let us now compute $S(\boldsymbol{\rho})$. In the case under consideration of a Coulomb field the integral in the exponent diverges, in correspondence with the divergence of the phase in Coulomb scattering. Therefore the integral must be taken between finite limits:
\[
\int_{-R}^{R}U\,dz = -2\nu\int_0^R\frac{dz}{\sqrt{\rho^2+z^2}} = -2\nu[\ln(R+\sqrt{R^2+\rho^2}) - \ln\rho] \approx -2\nu\ln R + 2\nu\ln\rho
\]
($R \gg \rho$). The first, constant, term is inessential, so that
\begin{equation}
S(\boldsymbol{\rho}) = \exp(-2i\nu\ln\rho) = \rho^{-2i\nu}.
\label{eq:96.17}
\end{equation}

Substituting (96.9), (96.17) in (96.15) and integrating over the directions of the vector $\boldsymbol{\rho}$ in the $xy$ plane, we find
\begin{equation}
a(\mathbf{q}_\perp) \propto \nu\int_0^{\infty}\rho^{-2i\nu}K_1(\chi)J_1(q_\perp\rho)\,\rho\,d\rho,
\label{eq:96.18}
\end{equation}

\SourcePage{470}{475}

where $J_1$ is the Bessel function. The factors not containing $\nu = Z\alpha$ are not written out here.

We see that the dependence of the amplitude $a(\mathbf{q}_\perp)$ (and consequently also of the cross section (96.16)) on $\nu$ is contained in a separate factor. On the other hand, when $\nu \to 0$ the cross section must tend to its Born value. Therefore it is clear that the cross section will differ from the Born one only by a factor that is independent of the polarisation of the electron and does not affect the polarisation effects.

The integral (96.18) can be expressed through a hypergeometric function with the aid of the formula
\[
\int_0^{\infty}x^{-\lambda}K_1(ax)J_1(bx)\,x\,dx = \frac{b\,\Gamma(2-\lambda/2)\,\Gamma(1-\lambda/2)}{2^\lambda a^{3-\lambda}}\!\left(1+\frac{b^2}{a^2}\right)^{\!-1+\lambda/2}\!F\!\left(\frac{\lambda}{2},\,1-\frac{\lambda}{2},\,2,\,\frac{b^2}{a^2+b^2}\right).
\]
This gives
\begin{equation}
a(\mathbf{q}_\perp) \propto \nu(1-i\nu)\!\left(\frac{q}{2}\right)^{\!2i\nu}\!\Gamma^2(1-i\nu)\,F(i\nu,\,1-i\nu,\,2,\,z),
\label{eq:96.19}
\end{equation}
where
\begin{equation}
z = 1 - \frac{m^4\omega^2}{4q^2\varepsilon^2\varepsilon'^2}(1+\delta^2)^2, \quad \delta = \frac{\varepsilon\theta}{m};
\label{eq:96.20}
\end{equation}
here it is used that in region II (see (96.2)) the component of the vector $\mathbf{q}$ parallel to $\mathbf{p}$ is
\begin{equation}
q_\parallel^2 = q^2 - q_\perp^2 \approx \frac{m^4\omega^2}{4\varepsilon^2\varepsilon'^2}(1+\delta^2)^2.
\label{eq:96.21}
\end{equation}

In this one can easily convince oneself, if one takes into account that in the indicated region the angles between the momenta $\mathbf{p}$, $\mathbf{p}'$ and $\mathbf{k}$ satisfy conditions (93.15).

The hypergeometric function in (96.19) can be reduced to the function $F(z)$ (95.15) with the aid of the formula
\[
F(a,\,b+1,\,c+1,\,z) = \frac{c}{c-a}\,F(a,\,b,\,c,\,z) + \frac{c(1-z)}{b(a-c)}\,F'(a,\,b,\,c,\,z).
\]

The final result takes then the form
\begin{equation}
d\sigma = d\sigma_\text{B}\,\frac{1}{F^2(1)}\!\left[F^2(z) + \frac{(1-z)^2}{\nu^2}\,F'^2(z)\right],
\label{eq:96.22}
\end{equation}
where $d\sigma_\text{B}$ is the Born cross section (93.13) (\emph{H.~A.~Bethe, L.~Maximon}, 1954). When $q \gg m^2/\varepsilon$ we have $z \approx 1$, so that the entire coefficient

\SourcePage{471}{476}

at $d\sigma_\text{B}$ tends to unity. In this sense formula (96.22), derived for region II, is automatically valid for all $q \lesssim m$. When $q \lesssim m^2/\varepsilon$ and the correction factor in (96.22) differs from unity, the vectors $\mathbf{p}$, $\mathbf{p}'$, $\mathbf{k}$ are almost coplanar and the quantities $\delta$ and $\delta'$ are almost equal to each other; this has already been taken into account in (96.22). Thus, $q$ in expression (96.20) for $z$ can be rewritten as
\begin{equation}
\frac{q^2}{m^2} = \delta^2 + \delta'^2 - 2\delta\delta'\cos\varphi + \frac{m^2\omega^2}{4\varepsilon^2\varepsilon'^2}(1+\delta^2)^2,
\label{eq:96.23}
\end{equation}
i.e.\ one can set $\delta = \delta'$ in the second term in (93.14), but not in the first term, which does not contain the small coefficient ($\sim m^2/\varepsilon^2$).

For finding the integral (over angles) cross section of radiation there is no need to carry out the integration anew, as is clear from the following considerations (\emph{H.~Olsen}, 1955). The different directions of $\mathbf{p}'$ (at given energy $\varepsilon'$) correspond to the degeneracy of the final state of the electron. It is evident that the result of summation over states belonging to one degenerate level does not depend on the manner in which the complete set of these states is chosen. We can therefore use for the purposes of summation over directions of $\mathbf{p}'$ the system of functions $\psi_{\varepsilon'\mathbf{p}'}^{(-)}$ instead of the system $\psi_{\varepsilon'\mathbf{p}'}^{(+)}$ (necessary for computing the differential cross section), i.e.\ define the matrix element of bremsstrahlung as
\[
M_{fi}^\text{brem} = \int\psi_{\varepsilon'\mathbf{p}'}^{(+)*}(\boldsymbol{\alpha}\mathbf{e}^*)\,e^{-i\mathbf{k}\mathbf{r}}\,\psi_{\varepsilon\mathbf{p}}^{(+)}\,d^3x.
\]

It is easy to verify that this integral coincides with the integral $(M_{fi}^\text{pair})^*$, if in the latter one replaces the parameters of the wave functions according to
\[
\mathbf{p}_+, p_+, \varepsilon_+ \to -\mathbf{p}, -p, -\varepsilon; \quad \mathbf{p}_-, p_-, \varepsilon_- \to \mathbf{p}', p', \varepsilon'; \quad \mathbf{k} \to -\mathbf{k}
\]
(and also replaces the variables of integration: $\mathbf{r} \to -\mathbf{r}$).

From this it is clear that the integral (over angles) bremsstrahlung cross section can be obtained from the integral pair-production cross section (95.20) by multiplying the latter by
\[
\frac{\omega^2\,d\omega}{p_+^2}\,\frac{d\omega}{\omega} \approx \frac{\omega^2\,d\omega}{\varepsilon_+^2\varepsilon_-}
\]
(cf.\ (91.6)) and replacing $\varepsilon_+ \to -\varepsilon$, $\varepsilon_- \to \varepsilon'$. Thus, we find
\begin{equation}
d\sigma = 4Z^2\alpha r_\mathrm{e}^2\,\frac{\varepsilon'}{\varepsilon}\!\left[\frac{\varepsilon}{\varepsilon'} + \frac{\varepsilon'}{\varepsilon} - \frac{2}{3}\right]\!\left(\ln\frac{2\varepsilon\varepsilon'}{m\omega} - \frac{1}{2} - f(\alpha Z)\right)\frac{d\omega}{\omega}.
\label{eq:96.24}
\end{equation}
